% ============================================================
% 片段：s10-refs  参考文献（原文 91 条，编号与原文一致）
% ============================================================

\clearpage
\begin{thebibliography}{99}

\bibitem{ref1} R. Langlands,
``Problems In The Theory Of Automorphic Forms,'' in Lect. Notes in
Math. {\bf 170} (Springer-Verlag, 1970), pp. 18-61; ``Where Stands
Functoriality Today?'' in {\it Representation Theory And
Automorphic Forms}, Proc. Symp. Pure Math. {\bf 61} (American
Mathematical Society, 1997), pp. 457-471.

\bibitem{ref2} A. Beilinson
and V. Drinfeld, ``Quantization Of Hitchin's Integrable System And
Hecke Eigensheaves,'' preprint (ca. 1995),
http://www.math.uchicago.edu/~arinkin/langlands/.

\bibitem{ref3} E. Frenkel, ``Affine Algebras, Langlands Duality,
And Bethe Ansatz,'' in the Proceedings of the XIth International
Congress of Mathematical Physics, Paris, 1994, ed. D. Iagolnitzer
(International Press, 1995) 606-642, q-alg/9506003.

\bibitem{ref4} A.
Beilinson and V. Drinfeld, {\it Chiral Algebras}, American
Mathematical Society Colloquium Publications {\bf 51}
(American Mathematical Society, 2004).

\bibitem{ref5} E. Frenkel
and D. Ben-Zvi, {\it Vertex Algebras And Algebraic Curves},
Mathematical Surveys And Monographs {\bf 88},
second edition (American Mathematical Society, 2004).

\bibitem{ref6} E. Frenkel and D. Gaitsgory, ``Local Geometric
Langlands Correspondence And Affine Kac-Moody Algebras,''
math.RT/0508382.

\bibitem{ref7} E. Frenkel, ``Ramifications Of The Geometric Langlands Program,''
math.QA/0611294.

\bibitem{ref8} A. Kapustin and E. Witten, ``Electric-Magnetic Duality
And The Geometric Langlands Program,'' hep-th/0604151.

\bibitem{ref9} E. Frenkel, ``Lectures On The Langlands Program And
Conformal Field Theory,'' arXiv:hep-th/0512172.

\bibitem{ref10} R. M. Rohm, {\it Some Current Problems In
Physics Beyond The Standard Model}, Princeton University Ph.D. Thesis (1985), unpublished.

\bibitem{ref11} P. B. Kronheimer and T. S. Mrowka,
``Gauge Theory For Embedded Surfaces, I, II,'' Topology {\bf 32}
(1993), 773-826, {\bf 34} (1995) 37-97.

\bibitem{ref12} E. Witten, ``Topological Sigma
Models,'' Commun. Math. Phys. {\bf 118} (1988) 411.

\bibitem{ref13} J. Preskill and
L. M. Krauss, ``Local Discrete Symmetry
And Quantum Mechanical Hair,''
Nucl. Phys. {\bf B341} (1990) 50-100.

\bibitem{ref14} M. G. Alford, K.-M. Lee, J. March-Russell, and J. Preskill, ``Quantum Field Theory
Of Non-Abelian Strings And Vortices,'' Nucl. Phys. {\bf B384} (1992) 251-317, hep-th/9112038.

\bibitem{ref15} M. Bucher, K.-M. Lee, and J. Preskill, ``On Detecting Discrete Cheshire Charge,''
Nucl. Phys. {\bf B386} (1992) 27-42, hep-th/9112040.

\bibitem{ref16} A. Braverman, ``Instanton Counting Via Affine Lie Algebras I:
Equivariant J-Functions of (Affine) Flag Manifolds and Whittaker Vectors,'' math.AG/0401409.

\bibitem{ref17} A. Braverman and P. Etinghof, ``Instanton Counting Via Affine Lie Algebras II:
from Whittaker Vectors to the Seiberg-Witten Prepotential,'' math.AG/0409041.

\bibitem{ref18} N. R. Constable, J. Erdmenger, Z. Guralnik, and
I. Kirsch, ``Intersecting $D3$-Branes And Holography,''
hep-th/0211222.

\bibitem{ref19} M. Bershadsky, A. Johansen, V.
Sadov, and C. Vafa, ``Topological Reduction Of $4-D$ SYM  To $2-D$
Sigma Models,'' Nucl. Phys. {\bf B448} (1995) 166-186,
arXiv:hep-th/9501096.

\bibitem{ref20} J. A. Harvey, G. W. Moore, and A. Strominger, ``Reducing
$S$ Duality To $T$ Duality,'' Phys. Rev. {\bf D52} (1995) 7161-7167.

\bibitem{ref21} N. Hitchin, ``The Self-Duality Equations On
A Riemann Surface,'' Proc. London Math. Soc. (3) {\bf 55} (1987)
59-126.

\bibitem{ref22} C. Simpson,
``Harmonic Bundles On Noncompact Curves,'' J. Am. Math. Soc. {\bf
3} (1990) 713-770.

\bibitem{ref23} H. Konno, ``Construction Of The Moduli Space Of
Stable Parabolic Higgs Bundles On A Riemann Surface,'' J. Math.
Soc. Japan {\bf 45} (1993) 253-276.

\bibitem{ref24} H. Nakajima, ``Hyper-Kahler Structures On
Moduli Spaces Of Parabolic Higgs Bundles On Riemann Surfaces,''
in {\it Moduli Of Vector Bundles}, Lecture Notes in Pure and Appl. Math
{\bf 179} (Marcel Dekker, New York, 1996), ed. M. Maruyama.

\bibitem{ref25} O. Biquard and P. Boalch, ``Wild Non-Abelian Hodge Theory On Curves,''
Compos. Math. {\bf 140} (2004) 179-204.

\bibitem{ref26} T. Springer,  ``A
Construction Of Representations Of Weyl Groups,''
Invent. Math. {\bf 44} (1978) 279-293.

\bibitem{ref27} D. Kazhdan and G. Lusztig, ``Proof Of The Deligne-Langlands Conjecture For Hecke
Algebras,'' Invent. Math. {\bf 87} (1987) 153-215.

\bibitem{ref28} D. Kazhdan and G. Lusztig, ``Equivariant $K$-Theory And Representations Of Hecke
Algebras, II'' Invent. Math. {\bf 80} (1985) 209-231.

\bibitem{ref29} R.
Bezrukavnikov, ``Noncommutative Counterparts Of The Springer Resolution,''
math.RT/0604445.

\bibitem{ref30} N.
Chriss and V. Ginzburg, {\it Representation Theory And Complex Geometry} (Birkhauser, Boston,
1997).

\bibitem{ref31} M. F. Atiyah and R. Bielawski, ``Nahm's Equations, Configuration
Spaces, and Flag Manifolds,'' Bull. Braz. Math. Soc. (NS) {\bf 33} (2002) 157-76,
math.RT/0110112.

\bibitem{ref32} D. Ben-Zvi and D. Nadler, to
appear.

\bibitem{ref33} E. Witten, to appear.

\bibitem{ref34} P. Kronheimer, ``A Hyper-Kahlerian Structure
On Coadjoint Orbits Of A Semisimple Complex Group,'' J. London Math. Soc. {\bf 42} (1990)
193-208.

\bibitem{ref35} P. Kronheimer, ``Instantons And The Geometry of The Nilpotent Variety,''
J. Diff. Geom. {\bf 32} (1990) 473-490.

\bibitem{ref36} O. Biquard, ``Sur les \'Equations de Nahm et la Structure de Poisson des Alg\'ebres
de Lie Semi-Simples Complexes,'' Math. Ann. {\bf 304} (1996) 253.

\bibitem{ref37} A. G. Kovalev, ``Nahm's Equations And Complex Adjoint Orbits,'' Quart. J. Math.
Oxford Ser. (2) {\bf 47} (1996) 41-58.

\bibitem{ref38} A. Kapustin,
``Wilson-'t Hooft Operators In Four-Dimensional Gauge Theories
And $S$-Duality,'' hep-th/0501015.

\bibitem{ref39} P. Goddard, J. Nuyts, and D. I. Olive, ``Gauge
Theories And Magnetic Charge,'' Nucl. Phys. {\bf B125} (1977)
1-28.

\bibitem{ref40} R.
Bott and L. Tu, {\it Differential Forms In Algebraic Topology}
(Springer, New York, 1982).

\bibitem{ref41} N.~Dorey, C.~Fraser, T.~J.~Hollowood and
M.~A.~C.~Kneipp, ``$S$-Duality in ${\EUN}=4$ Supersymmetric Gauge
Theories With Arbitrary Gauge Group,'' Phys.\ Lett.\ B {\bf 383},
422-428 (1996), arXiv:hep-th/9605069.

\bibitem{ref42} P. Argyres, A. Kapustin, N. Seiberg, ``On $S$-Duality For Non-Simply-Laced Gauge
Groups,'' arXiv:hep-th/0603048.

\bibitem{ref43} E. Witten, ``On $S$-Duality In Abelian Gauge
Theory,'' Selecta Math. {\bf 1} (1995) 383, hep-th/9505186.

\bibitem{ref44} T. H. Buscher, ``Path Integral
Derivation Of Quantum Duality In Nonlinear Sigma
Models,'' Phys. Lett. {\bf 201B} (1988) 466.

\bibitem{ref45} M. Rocek and E. Verlinde, ``Duality, Quotients,
And Currents,'' Nucl. Phys. {\bf B373} (1992) 630,
hep-th/9110053.

\bibitem{ref46} E. Witten, ``Dyons Of Electric Charge
$e\theta/2\pi$,'' Phys. Lett. {\bf B86} (1979) 283-7.

\bibitem{ref47} G. 't Hooft, ``Magnetic Monopoles In Unified
Gauge Theories,'' Nucl. Phys. {\bf B79} (1974) 276-284.

\bibitem{ref48} C. Vafa and E. Witten, ``A Strong
Coupling Test Of $S$-Duality,'' Nucl. Phys. {\bf B431} (1994) 3-77, hep-th/9408074.

\bibitem{ref49} N. J.
Hitchin, A. Karlhede, U. Lindstrom, and M. Rocek, ``Hyperkahler
Metrics and Supersymmetry,'' Commun. Math. Phys. {\bf 108} (1987)
535-589.

\bibitem{ref50} V. B. Mehta and C. S. Seshadri, ``Moduli Of
Vector Bundles On Curves With Parabolic Structures,'' Math. Ann. {\bf 248} (1980) 205-39.

\bibitem{ref51} M. S. Narasimhan and C. S. Seshadri, ``Stable And Unitary Vector Bundles
On A Compact Riemann Surface,'' Ann. Math. (2) {\bf 82} (1965) 540-67.

\bibitem{ref52} M. F. Atiyah and R. Bott, ``The Yang-Mills Equations Over
Riemann Surfaces,'' Phil. Trans. Roy. Soc. London {\bf A308}
(1982) 523-615.

\bibitem{ref53} S. Kumar and
M. S. Narasimhan, ``Picard Group Of The Moduli Space Of $G$-Bundles,'' Math. Ann.
{\bf 308} (1997) 155-173.

\bibitem{ref54} Y. Laszlo and C. Sorger, ``The Line Bundles On The Moduli Of
Parabolic $G$-Bundles Over Curves And Their Sections,'' alg-geom/9507002.

\bibitem{ref55} C. Pauly, ``Fibr\'es Parabolique de Rang 2 et Fonctions Th\^eta
G\'eneralis\'ees,'' Math. Z. {\bf 128} (1998) 31-50.

\bibitem{ref56} L. Jeffrey, ``The Verlinde Formula For Parabolic Bundles,'' math.AG/0003150.

\bibitem{ref57} S. Donaldson, ``Gluing Techniques In The Cohomology Of Moduli Spaces,'' in {\it Topological
Methods In Modern Mathematics} (Publish or Perish Press, Houston, 1993), pp. 137-170.

\bibitem{ref58} L.
 Jeffrey, ``Extended Moduli Spaces Of Flat Connections On Riemann Surfaces,'' Math. Annalen
 {\bf 298} (1994) 667-692.

\bibitem{ref59} E. Witten,
``Quantum Field Theory And The Jones Polynomial,'' Commun. Math. Phys.
{\bf 121} (1989) 351.

\bibitem{ref60} H. U. Boden and Y. Hu, ``Variations Of Moduli Of
Parabolic Bundles,'' Math. Annalen {\bf 301} (1995) 539-559.

\bibitem{ref61} H. U. Boden and K. Yokogawa,
``Moduli Spaces Of Parabolic Higgs Bundles And Parabolic $K(D)$ Pairs,'' Internat. J.
Math. {\bf 7} (1996) 573-598.

\bibitem{ref62} P.
Kronheimer, ``The Construction Of ALE Spaces As Hyper-Kahler Quotients,''
J. Diff. Geom. {\bf 29} (1989) 665-83.

\bibitem{ref63} R.
Bielawski, ``Lie Groups, Nahm's Equations, And Hyper-Kahler Manifolds,'' math.DG/0509515.

\bibitem{ref64} W. Nahm, ``The Construction Of All
Self-Dual Multimonopoles By The ADHM Method,'' in {\it Monopoles And Quantum Field Theory},
ed. N. S. Craigie et. al. (World Scientific, 1982) 87-95.

\bibitem{ref65} D. Kazhdan and G. Lusztig, ``A Topological Approach
To Springer's Representations,'' Adv. Math. {\bf 38} (1980) 222-228.

\bibitem{ref66} C. H. Taubes, ``Self-Dual Connections
On 4-Manifolds With Indefinite Intersection Matrix,'' J. Diff. Geom. {\bf 19} (1984)
517-560.

\bibitem{ref67} A. Beauville, ``Jacobiennes des Courbes
Spectral et Syst\`emes Hamiltoniens Compl\`etement Int\'grables,'' Acta Math. {\bf 164}
(1990) 211-35.

\bibitem{ref68} E. Markman and R. Donagi, ``Spectral Curves, Algebraically Completely
Integrable Hamiltonian Systems, And Moduli Of Bundles,'' in Lecture Notes in Math. vol. 1620
(Springer, 1996) pp.  1-119,  alg-geom/9507017.

\bibitem{ref69} M. Thaddeus and T. Hausel, ``Mirror Symmetry, Langlands Duality, And
The Hitchin System,'' math.AG/0205236.

\bibitem{ref70} A. Braverman and R. Bezrukavnikov, ``Geometric Langlands Correspondence
For ${\mathcal D}$-Modules In Prime Characteristic: The $GL(N)$ Case,'' math.AG/0602255.

\bibitem{ref71} R. Donagi and T. Pantev, ``Langlands Duality For
Hitchin Systems,'' math.AG/0604617.

\bibitem{ref72} A.
Kapustin and D. Orlov, ``Remarks On $A$-Branes, Mirror Symmetry,
and the Fukaya Category,'' J. Geom. Phys. {\bf 48} (2003) 84-99,
arXiv:hep-th/019098.

\bibitem{ref73} M. Kontsevich, ``Deformation Quantization Of
Algebraic Varieties,'' math.AG/0106006.

\bibitem{ref74} A. Beilinson and J. Bernstein, "A Proof of Jantzen Conjectures,"
  Advances in Soviet Mathematics 16, Part 1, pp. 150, AMS, 1993.

\bibitem{ref75} D. Gaitsgory, ``Construction Of Central Elements In
The Affine Hecke Algebra Via Nearby Cycles,'' math.AG/9912074.

\bibitem{ref76} S. Coleman, ``More
About The Massive Schwinger Model,'' Annals Phys. {\bf 101} (1976)
239.

\bibitem{ref77} N. Nitsure,
``Cohomology Of The Moduli Of Parabolic Vector Bundles,'' Proc. Indian Acad. Sci. Math. Sci.
{\bf 95} (1986) 61-77.

\bibitem{ref78} A. Kapustin and S. Sethi,  ``The Higgs Branch Of Impurity Theories,''
Adv. Theor. Math. Phys. {\bf 2} (1998) 571-591,
hep-th/9804027.

\bibitem{ref79} O. DeWolfe, D.
Z. Freedman, and H. Ooguri, ``Holography And Defect Conformal Field Theories,'' Phys. Rev. {\bf
D66} (2002) 025009, hep-th/0111135.

\bibitem{ref80} M. Douglas and G. Moore, ``$D$-Branes, Quivers, And ALE Instantons,''
hep-th/9603167.

\bibitem{ref81} E. Witten, ``Some Comments On String Dynamics,'' in I. Bars et. al., eds. (1996),
{\it Future Perspectives In String Theory},
hep-th/9507121.

\bibitem{ref82} P. Z. Kobzk and A. F.
Swann, ``Classical Nilpotent Orbits As Hyper-Kahler Quotients,'' Int. J. Math. {\bf 7} (1996)
193-210.

\bibitem{ref83} L. Dixon, J. Harvey, C. Vafa, and E. Witten,
``Strings On Orbifolds,'' Nucl. Phys. {\bf B261} (1985) 678-686.

\bibitem{ref84} P. Aspinwall, ``Enhanced Gauge Symmetries And
K3 Surfaces,'' hep-th/9507012.

\bibitem{ref85} J. Maldacena, ``The Large $N$ Limit Of
Superconformal Field Theories And Supergravity,'' Adv. Theor. Math. Phys. {\bf 2} (1998) 231-252,
hep-th/9711200.

\bibitem{ref86} J. Maldacena, ``Wilson Loops In Large $N$ Field
Theories,'' Phys. Rev. Lett. {\bf 80}
(1998) 4859-4862, hep-th/9803002.

\bibitem{ref87} S.-J. Rey and J.-T. Yee, ``Macroscopic Strings As Heavy Quarks
in Large $N$ Gauge Theory And Anti-de Sitter Supergravity,'' Eur. Phys. J. {\bf C22} (2001) 379-394,
hep-th/9803001.

\bibitem{ref88} D. Berenstein, R. Corrado, W. Fischler, and J. M. Maldacena, ``The Operator
Product Expansion For Wilson Loops And Surfaces In The Large $N$ Limit,'' Phys. Rev.
{\bf D59} (1999) 105023, hep-th/9809188.

\bibitem{ref89} C. R. Graham and E. Witten, ``Conformal Invariance Of Submanifold Observables In
AdS-CFT Correspondence,'' Nucl. Phys. {\bf B546} (1999) 52-64, hep-th/9901021.

\bibitem{ref90} J. E.
Humphreys, {\it Introduction To Lie Algebras And Representation Theory} (Springer-Verlag, New York,
1972).

\bibitem{ref91} B. Simon, {\it Representations Of Finite And Compact Groups}
(American Mathematical Society, 1996).

\end{thebibliography}
